%! Skills Focus: Selecting an Appropriate Inference Procedure for Categorical Data — Unit 8, Lesson 8.7 — Exit Ticket (Answer Key)
\documentclass[10pt]{article}
\usepackage{apstats-article}
\usepackage{apstats-key}   % pulls in -boxes; do NOT also load -boxes

\begin{document}

\pageheader{Unit 8, Lesson 8.7}{Exit Ticket --- Answer Key}
\namedateperiod

\vspace{0.12in}

\textit{For each scenario: (a) name the correct inference procedure, and (b) write $H_0$ in
the appropriate parameter language. No computation is required.}

\vspace{0.10in}
\begin{enumerate}

\item A researcher collects one random sample of 300 college students and records two
      variables: each student's political affiliation (Democrat, Republican, Independent, Other)
      and whether the student voted in the last election (yes or no). The researcher wants to
      test for an association between political affiliation and voter participation.

      \begin{enumerate}[label=\alph*.]
        \item Procedure: \ans{chi-square test for independence}
        \item $H_0$ (write in context): \ansline{There is no association between political affiliation and voter participation among college students.}
      \end{enumerate}

\vspace{0.14in}

\item A candy company claims its fruit snack bags are distributed as follows: 25\% strawberry,
      25\% grape, 25\% orange, and 25\% cherry. A quality inspector randomly samples 160 bags and
      counts the flavor of one piece in each. He wants to test whether the actual distribution
      matches the company's claim.

      \begin{enumerate}[label=\alph*.]
        \item Procedure: \ans{chi-square goodness-of-fit test}
        \item $H_0$ (write in context): \ansline{The distribution of flavors in fruit snack bags is 25\% strawberry, 25\% grape, 25\% orange, and 25\% cherry.}
      \end{enumerate}

\vspace{0.14in}

\item A public health researcher randomly samples 200 adults from a rural county and 200 adults
      from an urban county. She records whether each adult has health insurance (yes or no). She
      wants to estimate how much higher the insured proportion is in the urban county.

      \begin{enumerate}[label=\alph*.]
        \item Procedure: \ans{two-sample $z$-interval for $p_1 - p_2$}
        \item Why is a confidence interval (not a test) appropriate here? \ansline{The goal is to \emph{estimate} the magnitude of the difference in proportions (``how much higher''), not to test a specific claim about equality. A CI provides a range for $p_{\text{urban}} - p_{\text{rural}}$.}
      \end{enumerate}

\end{enumerate}

\end{document}
